\documentclass[11pt]{article}

% ======================================================================
% OrbitalWatch: An Open-Source Framework for Detecting Anomalous Satellite
% Behavior and Cyber-Physical Risk Signals from Public Orbital Data
% Author: Khadija (Kadi) Taki
% Draft status: STRUCTURED DRAFT. Framing and methods are written.
% Results/tables/figures are marked [PENDING ANALYSIS] and must not be
% filled with invented numbers. Populate after you run the analysis.
% ======================================================================

\usepackage[utf8]{inputenc}
\usepackage[T1]{fontenc}
\usepackage[margin=1in]{geometry}
\usepackage{amsmath,amssymb}
\usepackage{booktabs}
\usepackage{graphicx}
\usepackage{xcolor}
\usepackage{hyperref}
\usepackage{enumitem}
\usepackage{authblk}

\hypersetup{colorlinks=true, linkcolor=black, citecolor=black, urlcolor=blue}

\newcommand{\pending}[1]{\textcolor{red}{\textbf{[PENDING ANALYSIS: #1]}}}
\newcommand{\fill}[1]{\textcolor{orange}{\textbf{[FILL: #1]}}}

\title{OrbitalWatch: An Open-Source Framework for Detecting Anomalous\\
Satellite Behavior and Cyber-Physical Risk Signals from Public Orbital Data}
\author[1]{Khadija Taki}
\affil[1]{Heinz College of Information Systems and Public Policy, Carnegie Mellon University\\ \texttt{\fill{preferred contact email}}}
\date{\today}

\begin{document}
\maketitle

\begin{abstract}
Public orbital tracking data, distributed as two-line element sets (TLEs),
provides a continuously updated, openly available record of the motion of
thousands of satellites. This paper asks whether that public data can be used to
detect statistically unusual satellite behavior that warrants further
cyber-physical investigation. We present OrbitalWatch, an open-source framework
that ingests public TLE data, derives per-satellite behavioral features
(inclination, eccentricity, mean motion, altitude, maneuver frequency, and
proximity to other objects), and applies a layered set of statistical anomaly
detection methods to flag deviations from each satellite's own historical
baseline and from its peer category. We are explicit about what such a signal
does and does not mean: an orbital anomaly is a cyber-physical risk signal
requiring corroboration, not proof of a cyberattack. We describe the data
pipeline, the detection methods, and a sensitivity and false-positive analysis
protocol, and we release the dataset, code, and visualizations to support
reproduction. \pending{quantitative findings, once the analysis is run}
\end{abstract}

\section{Introduction}

The space domain is now critical infrastructure. Communications, navigation,
Earth observation, and defense all depend on satellites whose behavior is, in
principle, observable from the ground. Two-line element sets, published for a
large fraction of tracked objects, encode each satellite's orbit compactly
enough to be analyzed at scale. If meaningful behavioral anomalies can be
extracted from this public record, then open data becomes a usable input to
space situational awareness and to the governance of space-based systems.

The central research question of this paper is deliberately narrow and testable:

\begin{quote}
\emph{Can public orbital data be used to detect statistically unusual satellite
behavior that warrants further cyber-physical investigation?}
\end{quote}

The framing matters. We are not claiming that an anomaly proves hostile action.
Orbital anomalies have many benign explanations: routine station-keeping,
collision avoidance, end-of-life disposal, or noise in the underlying tracking
data. Our claim is weaker and more defensible: that certain deviations, relative
to a satellite's own history and to its peer category, are statistically unusual
enough to justify corroboration from other sources. That methodological
restraint is what makes the work credible.

\paragraph{Contributions.}
\begin{enumerate}[leftmargin=*]
  \item An open pipeline that turns public TLE data into per-satellite
        behavioral time series (Section~\ref{sec:data}).
  \item A layered anomaly-detection methodology combining classical statistical
        baselines with an unsupervised detector (Section~\ref{sec:methods}).
  \item A sensitivity and false-positive analysis protocol appropriate to a
        domain with no ground-truth labels (Section~\ref{sec:eval}).
  \item A released dataset, code, and set of visualizations enabling
        reproduction (Section~\ref{sec:artifacts}).
\end{enumerate}

\section{Related Work}
\fill{2 to 3 paragraphs situating this against space situational awareness,
prior TLE-based analyses, and anomaly detection in astrodynamics. Add real
references to the .bib. Note honestly where OrbitalWatch differs: openness,
reproducibility, and the cyber-physical framing rather than pure astrodynamics.}

\section{Data}
\label{sec:data}

\subsection{Source and scope}
We collect public TLE data for \pending{number} satellites over an observation
window of \pending{start date to end date}. Satellites are grouped into
functional categories, communications, navigation, Earth observation, and
defense-related, so that anomalies can be judged both against a satellite's own
history and against its category peers.

\subsection{Derived features}
From successive TLEs we derive the following per-satellite features:
\begin{itemize}[leftmargin=*]
  \item Changes in inclination
  \item Changes in eccentricity
  \item Changes in mean motion
  \item Unexpected altitude changes
  \item Unusual proximity between space objects
  \item Maneuver frequency
  \item Deviations from each satellite's historical behavior
  \item Differences across the communications, navigation, Earth-observation,
        and defense-related categories
\end{itemize}

\subsection{Preprocessing}
\fill{Describe cleaning: handling of missing or duplicate TLEs, epoch alignment,
outlier handling in the raw elements, and any smoothing. State the final
per-satellite series length.}

\section{Methods}
\label{sec:methods}

OrbitalWatch layers several detectors so that no single method carries the
finding alone. Simple, interpretable baselines run first; an unsupervised
detector adds sensitivity to multivariate patterns.

\begin{itemize}[leftmargin=*]
  \item \textbf{Rolling averages and standard deviations} establish each
        satellite's own moving baseline.
  \item \textbf{Robust Z-scores} flag deviations while resisting the influence
        of extreme single points.
  \item \textbf{Median absolute deviation (MAD)} provides a robust dispersion
        measure for skewed or heavy-tailed features.
  \item \textbf{Change-point detection} identifies shifts in a satellite's
        behavioral regime rather than isolated spikes.
  \item \textbf{Isolation Forest} provides unsupervised multivariate anomaly
        scoring across the derived features~\cite{iforest}.
  \item \textbf{Time-series visualization} supports human review of any flagged
        satellite.
\end{itemize}

\fill{For each method, state the exact parameters used (window sizes, Z-score and
MAD thresholds, change-point penalty, Isolation Forest contamination and
estimators). Reproducibility depends on these being written down.}

\section{Evaluation Protocol}
\label{sec:eval}

Because there is no public ground truth for ``a satellite that was actually
interfered with,'' we cannot report conventional precision and recall against
labels. Instead we evaluate along axes appropriate to an unlabeled,
high-consequence domain:

\begin{itemize}[leftmargin=*]
  \item \textbf{Sensitivity analysis:} how flagged counts change as thresholds
        vary, reported as a curve rather than a single operating point.
  \item \textbf{False-positive analysis:} manual review of a sample of flags
        against known benign causes (station-keeping, disposal, avoidance).
  \item \textbf{Cross-method agreement:} whether independent detectors flag the
        same satellite-windows, as a proxy for robustness.
  \item \textbf{Category baselining:} whether flags concentrate in categories in
        a way consistent with known operational behavior.
\end{itemize}

\section{Results}
\pending{This entire section awaits your analysis run. Do not populate with
placeholder numbers. Suggested structure once results exist: (1) overall flag
counts and rates by method; (2) sensitivity curves; (3) cross-method agreement;
(4) a small number of illustrative flagged cases with the explicit caveat that
each requires corroboration.}

\begin{table}[h]
\centering
\caption{Anomaly flags by detector and category. \pending{populate}}
\begin{tabular}{@{}lccc@{}}
\toprule
\textbf{Category} & \textbf{Robust Z} & \textbf{Change-point} & \textbf{Isolation Forest} \\
\midrule
Communications & \pending{} & \pending{} & \pending{} \\
Navigation & \pending{} & \pending{} & \pending{} \\
Earth observation & \pending{} & \pending{} & \pending{} \\
Defense-related & \pending{} & \pending{} & \pending{} \\
\bottomrule
\end{tabular}
\end{table}

\section{Discussion: Interpreting an Orbital Anomaly}

We restate the central methodological point because it governs how any result
here should be read. An orbital anomaly does not prove a cyberattack. It is a
cyber-physical risk signal requiring corroboration. A responsible use of
OrbitalWatch treats a flag as a prompt to gather additional evidence, not as an
attribution. This restraint is not a weakness of the framework; it is the
condition under which open-source orbital analysis can contribute to governance
without producing false accusations.

\section{Released Artifacts}
\label{sec:artifacts}
\fill{List the released dataset, the code repository URL, the figures, and the
two-page policy brief. A persistent link (for example a Zenodo DOI) strengthens
this considerably for a writing sample.}

\section{Limitations}
TLE data is itself imperfect: it is a fitted approximation, updated at irregular
cadence, and subject to its own errors. Absence of a flag does not mean absence
of activity, and presence of a flag has many benign explanations. The framework
detects statistical unusualness, not intent.

\section{Conclusion}
OrbitalWatch tests whether openly available orbital data can surface
statistically unusual satellite behavior worth further investigation, and
packages the result as reproducible open science. Its contribution to governance
is a disciplined, corroboration-first way to use public data as an input to
space situational awareness. \pending{one closing sentence stating the empirical
finding, once known}

\begin{thebibliography}{9}
\bibitem{iforest}
F.~T. Liu, K.~M. Ting, and Z.-H. Zhou,
``Isolation Forest,''
\emph{2008 Eighth IEEE International Conference on Data Mining}, 2008, pp.\ 413--422.
% \fill{Add: TLE / space-track data source citation, space situational awareness
% references, change-point and robust-statistics references you rely on.}
\end{thebibliography}

\end{document}
